% =====================================================================
\section{The sixth species: five classes closed, a quantitative data-mode
theorem, and the two constants that were wrong}
\label{ch:sixth}
% =====================================================================

The charter for this phase (\code{research/v3/SIXTH-SPECIES-CHARTER.md}) replaced the slogan
``find a new idea'' with a specification. A candidate must satisfy \textbf{five conditions}
simultaneously --- nonperturbative at $\beta\approx2.5$; a \emph{sign-controlling lower} bound;
non-circular (it may not presuppose the gap, the threshold, or the isolation); not SUSY and
\emph{failing} on the free massless field, where $f_1^2=0$ and $G''=0$; and uniform in $(L,a)$,
surviving $a\to0$ --- and it must be scored on a \textbf{five-item acid battery}: report zero on the
free field, refuse a certificate to a thin atom, fail on the deep double well, keep the total
normalization $\operatorname{Var}(q)=\chi$ distinct from the mean normalization
$\operatorname{Var}(\bar q)=\chi/V$, and carry the right dimension ($G''_a=a^2m_{\mathrm{phys}}^2$,
never the equipartition budget $E[f^2]\approx2$). The charter also names the three admissible
outcomes: $(\alpha)$ a fifth data mode, computable \emph{and} floor-sensitive; $(\beta)$ a new
equivalent coordinate whose positivity is the gap, equipped with a machine strictly stronger than
$H1$ somewhere; $(\gamma)$ a decisive refutation that tightens the map.

This chapter reports the first two waves run against that specification. The one-sentence result
is: \textbf{five classes are closed with quantifiable causes; the meta-theorem $F3\wedge F4=\varnothing$
acquires a quantitative form but only strictly inside the spectral class; the named inequality
$(NA)$ is false in the only form in which it was ever stated; the third named inequality
$(\mathrm{TCI})$ turns out to be under-specified rather than merely unproved; two load-bearing
constants are wrong; and two bridges were built that turn ``alignment'' and ``ratio'' into one
comparison each, with the same missing input.} No claim below is a proof of the gap. Clay remains
$0\%$.

\subsection{Five classes closed: the anatomy, in one table}

A \emph{species} here means a mechanism class, not a renamed inequality: if the class dies, every
future tool built out of it inherits the death. The five candidates of the first wave are scored in
Table~\ref{tab:species16}; the numbers in the ``decisive cause'' column are the ones that do the
killing, and they are reproduced from \code{ROUND1-VERDICTS.md} \S\S1--5 with the corrections
required by the independent audit \code{AUDIT-R1.md}. Each row's label is a class-level statement:
\LN{} for a no-go of the class, \LT{} for the pieces that are theorems.

\footnotesize
\begin{longtable}{@{}c>{\raggedright\arraybackslash}p{2.05cm}
  >{\raggedright\arraybackslash}p{4.35cm}>{\raggedright\arraybackslash}p{6.35cm}c@{}}
\caption{The five species of the first wave, their decisive cause with the number that carries it,
and the label. The right-hand column is the class-level verdict, not a verdict on a single script.}
\label{tab:species16}\\
\toprule
\textbf{\#} & \textbf{Species} & \textbf{What it tries} & \textbf{Decisive cause, with number} &
\textbf{Label}\\
\midrule
\endfirsthead
\toprule
\textbf{\#} & \textbf{Species} & \textbf{What it tries} & \textbf{Decisive cause} &
\textbf{Label}\\
\midrule
\endhead
\bottomrule
\endlastfoot
1 & Harish--Chandra spherical (HC) &
read the conditioned marginal's Stieltjes transform as a spherical transform; locate the atom as a
Plancherel singularity, a boundary mass, or a point of a compact quotient &
$\Phi=-\log J_{\mathrm{Haar}}=-\log\lvert c_{\mathrm{SU}(N)}/c_{\mathrm{flat}}\rvert^2$ (residual
$8.5\times10^{-16}$) is a definitional identity, not a bridge; split-rank Plancherel is absolutely
continuous $\Rightarrow$ no discrete series; the $H^{1/2}$ norm diverges logarithmically
($20.72\to36.84$); compact quotient $=$ pure Peter--Weyl, levels $(0,2,6,12)$ independent of
$\beta$ $\Rightarrow$ the free field passes too &
\LN\\
2 & Pseudospectral non-normality (PS) &
pseudospectral inclusion $\lVert(T-z)^{-1}\rVert\le1/\varepsilon$ plus Rouch\'e/argument principle
$\Rightarrow$ ``no spectrum in the contour'' $\Rightarrow$ $m\ge r>0$ (the correct direction) &
$\varepsilon_1=0.7871$ against $\operatorname{dist}(\text{contour},\sigma)=0.001944$, so
$\varepsilon_1/\operatorname{dist}=404.8$ at $V=1$; the formal condition is
$V<\operatorname{dist}/\varepsilon_1=2.47\times10^{-3}<1$, admitting no positive integer; the
abelian rotor has the same defect &
\LN\\
3 & Tropical min-plus (TR) &
$M=e^{\beta W}$, max-plus cycle mean $\rho_\infty$ plus a Collatz--Wielandt/log-sum-exp squeeze
$\Rightarrow$ a two-sided bound &
the constant-modulus slice has $S_W\equiv0$ exactly, so the critical graph is the whole diagonal
and $\Delta_{\mathrm{trop}}\equiv0$ for every $\beta$, against the exact $\Delta=0.6444$;
$\mathrm{width}/\Delta=3.013$ at $\beta=2.5$ ($2.415$ after blocking to $L=6$); Dobrushin
$-\log\delta=0.0697$ is weaker by $9.2\times$ &
\LN\\
4 & Lorentzian / M-convex (LX) &
label-space log-concavity and Lorentzian/M-convex order $\Rightarrow$ a TP$_2$-type isolation &
the frustrated plaquette's label weight is $(-1)^{2j}w_j$: $n_{\text{negative,twisted}}=0$ yet the
order does not preserve sign (partial sums converge to $0.5034$), so a sign-definite
M-convex-supported weight has no referent here; stretched LR log-concavity is itself unproved in
general &
\LN\\
5 & Setwise / small-ball (SW) &
$(D,\alpha,\rho''(\text{mode})\ge\kappa)\Rightarrow\operatorname{Var}\le C/\kappa$ from setwise
geometry &
$\mu\mapsto\lambda\mu$ preserves $(D,\alpha)$ while $\operatorname{Var}\mapsto\lambda^2
\operatorname{Var}$ ($\operatorname{Var}=1891.27$ at $\varepsilon=0.1$, $R=100$); tail placement
preserves all local data while $\operatorname{Var}\sim\varepsilon R^2$;
$(\log\rho)''=-1/(2x^2)-1/(2(2-x)^2)<0$ on all of $(0,2)$, so the hypothesis
$\rho''(\text{mode})\ge\kappa>0$ is empty &
\LN\\
\end{longtable}
\normalsize

\textbf{HC: what dies, and what does not.} The species wanted the atom to be a singularity of a
spherical (Harish--Chandra) transform. Two of its three readings die outright and the third dies on
the acid test. \emph{First}, on a split-rank noncompact symmetric space the Plancherel measure is
absolutely continuous: the density is smooth and everywhere positive on the positive axis (value
$2.5\times10^{-11}$ at $\lambda=1$, asymptotic slope $1/4\pi=0.0796$), and the discrete mass at the
floor is $0.0$ --- there is no discrete series to supply an atom \LN. \emph{Second}, the reading
that moves the atom to the boundary mass puts it outside $\mathcal H_{1/2}$: the truncated norm
grows $20.72\to36.84$ as the cutoff runs from $10^2$ to $10^8$, i.e.\ logarithmically, so the
``boundary mass'' is not an element of the space in which the spectral statement lives \LN.
The audit adds a provenance note that must be kept: the sequence $20.72\to36.84$ is the analytic
formula $\log(Y/10^{-8})$ evaluated at two cutoffs, not a quadrature, so its label is analytic and
the boundary-mass route has no valid numerical content behind it (\code{AUDIT-R1.md} \S D.2).
\emph{Third}, the compact-quotient reading is killed by the acid test in the sharpest possible way:
a compact quotient has pure point spectrum by Peter--Weyl, so \emph{everything} is an atom, and the
levels are the Casimir values $(0,2,6,12)$ independent of $\beta$. A statement that is true of the
free field at every coupling cannot be the gap.

The surviving positive content of the species is smaller than its first report claimed, and the
correction is instructive. The crosswalk
\[
  \Phi=-\log J_{\mathrm{Haar}}=-\log\lvert c_{\mathrm{SU}(N)}/c_{\mathrm{flat}}\rvert^2,
  \qquad c_{\mathrm{flat}}=\prod_{\alpha>0}\langle\alpha,X\rangle,
  \label{eq:hccrosswalk}
\]
with $c_{\mathrm{SU}(N)}$ the Weyl denominator, holds to a residual $8.5\times10^{-16}$ on
$\mathrm{SU}(3)$; but both sides are the same product, so \eqref{eq:hccrosswalk} is the
\emph{definition} of the Weyl denominator, not a bridge from lattice gauge theory to noncompact
harmonic analysis, and the numeric agreement is between two routines that share a basis bug, hence
not independent evidence (\code{AUDIT-R1.md} \S A1.3--A1.4, \S C-6). What is genuine is the
curvature floor: $\min\Phi''=1/6=N/12$ arises as the per-root sum of $g''(0)=1/6$ with
$g(t)=-2\log(\sin(t/2)/(t/2))$, and the audit confirms $\operatorname{Hess}\Phi|_0=(N/12)\delta_{ab}$
as an $\mathrm{Ad}$-invariant scalar with three independent algorithms (finite-difference Hessian,
Rayleigh quotient, closed Cartan form; spread $\le4.3\times10^{-9}$, values $1/6,1/4,1/3$ for
$N=2,3,4$) \LT. This is the Haar/Jacobian floor already used in the main theorem, not a new
identity; the first wave's claim of a new identity is withdrawn.

\textbf{PS: the class-level kill, and the correction that strengthens it.} The pseudospectral route
is the only candidate of the five whose \emph{direction} was right: a resolvent bound plus Rouch\'e
really does exclude spectrum from a contour, hence gives a lower bound. The inclusion lemma itself
is a theorem \LT. Three things then kill the class. (i) It is not a non-abelian mechanism: the
abelian rotor with kernel $e^{\beta\cos\varphi}$ on the circle already has
$\mathrm{dep}>0$ for $\beta>0$ \LN. (ii) The frame cost is fatal in both frames: the conditioning
number of the similarity that symmetrizes the operator is
$\mathrm{cond}(e^{-W/2})=e^{(W_{\max}-W_{\min})/2}$, which is bounded for the rotor ($12.18$ at
$\beta=2.5$, $54.6$ at $\beta=4$) but unbounded in the co-moving frame ($2.9\times10^{215}$ at
$R=8$, overflow at $R=10$); in the flat frame the resolvent bound is therefore vacuous, and in the
weighted frame the operator is \emph{exactly} self-adjoint ($\kappa=1$), so the non-normality is a
gauge artifact \LN. (iii) The decisive number is the one-body comparison. With the $\beta=0$
operator as reference, the normalized model distance is $\varepsilon_1=0.7871$ and the distance
from the contour to the spectrum is $0.001944$, so
\[
  \frac{\varepsilon_1}{\operatorname{dist}(\text{contour},\sigma)}=404.8\ \ (V=1),
  \qquad
  V<\frac{\operatorname{dist}}{\varepsilon_1}=2.47\times10^{-3}<1 .
  \label{eq:ps404}
\]
Rouch\'e needs a positive integer $V$ below the right-hand side; there is none. The one-body
problem already fails by a factor $404.8$, and the volume never enters \LN.

The audit's correction matters because it changes the diagnosis rather than the verdict. The tensor
estimate $\lVert T^{\otimes V}-T_0^{\otimes V}\rVert\lesssim V\varepsilon_1$ is a first-order
expansion, valid only while $V\varepsilon_1\ll1$ --- here $V\le2$; for $V\ge3$ the hard bound
$\lVert\cdot\rVert\le2$ truncates it, so the rows $V=8,16,256,65536$ of the original table are
vacuous extrapolations and the single decisive row is $V=1$ (\code{AUDIT-R1.md} \S B2). The audit's
diagonal counterexample makes the direction of the error explicit: for
$T_0=\mathrm{diag}(1,0.1)$, $T=\mathrm{diag}(1,0.8871)$ the true distance falls
$0.7871\to0$ as $V$ runs $1\to65536$, while $V\varepsilon_1$ rises to $5.2\times10^{4}$ \LE. And the
attribution is corrected as well: the death is \emph{not} the $e^{-cV}$ volume barrier of the
Gauss-invariant Perron--Frobenius route (which is exponentially small in volume) but the fact that
$T(2.5)$ sits an $O(1)$ model distance away from every solvable reference, compared with its gap.
The failure is a nonperturbative comparison failure, not a volume-scaling one (\code{AUDIT-R1.md}
\S C-3). The reusable yield is that the entire family ``count spectrum with a resolvent, a contour,
or a pseudospectrum'' is closed in one stroke.

\textbf{TR: the direction account and the width account both fail.} The tropical route takes
$M=e^{\beta W}$, reads the max-plus cycle mean $\rho_\infty$ (the $\beta\to\infty$ dequantization
of the dominant eigenvalue), and adds a Collatz--Wielandt or log-sum-exp squeeze to obtain a
two-sided bound. On the constant-modulus slice the Wilson action is exactly flat, $S_W\equiv0$
(each plaquette is a pure difference, so the action vanishes identically), and the critical graph
is the whole diagonal: all $n=320$ self-loops tie at $\rho_\infty$, the max-plus spectrum is
$n$-fold degenerate, and $\Delta_{\mathrm{trop}}\equiv0$ for \emph{every} $\beta$ --- against the
exact transfer gap $\Delta=0.6444$ of the periodic anchor \LN. The best $2$-cycle mean is strictly
lower, $2.487962<2.5$, the deficit being $\beta(1-\cos 2\pi/n)=O(\beta/n^2)\to0$, so this is
structural blindness, not ``almost right''.
The squeeze width then exceeds the signal: $\mathrm{width}/\Delta=3.013$ at $\beta=2.5$ (only at
$\beta=20$ does it fall to $0.94$), and it is resolution-dependent rather than convergent,
$1.718\to3.443$ for $n_z=40\to640$ (approximately $\log n/\beta$). Blocking, the campaign's
standard remedy, does not save it: $w_L/\Delta$ moves only $3.013$ ($L=1$) $\to2.415$ ($L=6$), with
$N_{\mathrm{eff}}^{(L)}=e^{Ls}$ and $s_L=4.85\to3.89>0$ \LN. Two further controls agree: in the
frustration-free flat model $\Delta_{\mathrm{trop}}=0$ against an exact $\Delta=0.267884$ that is
$n$-independent ($n=64$ and $n=4096$ give the same value), with $\mathrm{width}/\Delta=10.465$ at
$n_z=4096$ ($4.255$ at $n_z=64$); and
the bidiagonal deficit gives $\beta\delta_2=8.7\times10^{-4}$ against $\Delta=0.6444$, short by
$750\times$. The Dobrushin side is short by $9.2\times$ ($\delta(P)=0.9327$, hence
$-\log\delta=0.0697$).

The structural sentence this wave contributes is worth keeping, with the audit's scope attached.
The gap is carried \emph{entirely} by the prefactor --- the Van Vleck Gaussian width $1/\sqrt\beta$
times the compact measure --- while in the noncompact (free) case the same operator has continuous
spectrum and no gap. Hence \textbf{compactness is a property of the measure, not of the action, and
the tropical object sees only the action} \LT{} (the max-plus spectrum statement), while the
compactness reading is a structural identification of the measure, not a missing step \LC{}
(\code{AUDIT-R1.md} \S C-8). The audit's warning is accepted in full: the
theorem covers max-plus cycle-mean data, not ``all tropical objects''; a future construction that
carries the Van Vleck prefactor inside the tropical object is not excluded by this evidence.

\textbf{SW: three independent proofs that the hypothesis is empty.} The setwise route proposed
$(D,\alpha,\rho''(\text{mode})\ge\kappa)\Rightarrow\operatorname{Var}\le C/\kappa$, where $D$ is a
doubling constant, $\alpha$ a small-ball exponent, and $\rho''(\text{mode})$ the curvature of the
marginal at its mode. Three separate failures close it. (i) \emph{Scale invariance}: the map
$\mu\mapsto\lambda\mu$ leaves $(D,\alpha)$ unchanged and multiplies the variance by $\lambda^2$;
the witness is $\varepsilon=0.1$, $R=100$, $\operatorname{Var}=1891.27$ \LN. (ii) \emph{Tail
placement}: $\mu\mapsto(1-\varepsilon)\mu+\varepsilon\delta_R$ leaves every local datum unchanged
while $\operatorname{Var}\sim\varepsilon R^2$ \LN. (iii) \emph{The hypothesis is empty}: the exact
one-plaquette marginal is $\rho_\beta(x)\propto e^{-\beta x}\sqrt{x(2-x)}$ on $x=1-\tfrac12
\operatorname{tr}U\in[0,2]$, and
\[
  (\log\rho)''(x)=-\frac{1}{2x^{2}}-\frac{1}{2(2-x)^{2}}<0
  \qquad\text{for all }x\in(0,2),
  \label{eq:swlogcurv}
\]
with value $-1$ at $x=1$. The audit requires two separate sentences here, and both are true
(\code{AUDIT-R1.md} \S B3, \S C-4). \emph{At fixed $x$}, the right-hand side of \eqref{eq:swlogcurv}
does not involve $\beta$ at all, because $-\beta x$ is linear and drops out of the second
derivative; a finite-difference check at $x=0.2,1.0,1.8$ for $\beta=0,2.5,4$ agrees to
$\le2.7\times10^{-7}$. \emph{At the mode}, the value does move with $\beta$: the mode solves
$(1-x)/(x(2-x))=\beta$, i.e.\ $x_*(\beta)=[(2\beta+1)-\sqrt{4\beta^2+1}]/(2\beta)$, and
$(\log\rho)''(x_*)$ runs $-1.0000$ ($\beta=0$), $-1.7071$ ($0.5$), $-3.6180$ ($1.0$), $-15.5495$
($2.5$), $-36.5311$ ($4.0$) \LE. The two sentences are consistent because the evaluation point
moves; the composite is negative for every $\beta$, since $\rho>0$ and $\rho'(x_*)=0$ give
$\rho''(x_*)=\rho(x_*)(\log\rho)''(x_*)<0$. One caveat is preserved from the audit: $\rho''$ itself
is \emph{not} negative on all of $(0,2)$ at large $\beta$ ($\beta=10$, $x=0.5$ gives
$\rho''/\rho=84.89>0$); only $(\log\rho)''$ is everywhere negative, and the text now says so.
(iv) The free Gaussian has $(D,\alpha)$ independent of $\sigma$ while $\operatorname{Var}=\sigma^2$,
so the acid test has nothing to bite on \LN.

The verdict on SW is the sharpest of the five because it inverts the promise of the class: the data
are genuinely computable ($\nu(A)$ is a finite datum), but they are floor-blind \emph{and}
direction-reversed --- small-ball bounds say ``mass cannot concentrate'', while a gap says ``mass
does concentrate at $m_1^2>0$''. The independent literature check reaches the same conclusion from
the analytic side (\code{LIT-CHECK.md} Q2-5), which is why the inversion is not an artifact of the
battery.

\textbf{LX: not a false hypothesis but a missing referent.} The Lorentzian route needs a
sign-definite, M-convex-supported label weight. On $\mathrm{SU}(2)$ the label weight
$w_j=(2j+1)c_j(s)$ is positive and log-concave at $\beta=2.5$, but it is only the fusion-graph
label weight, i.e.\ the same object that already died as \code{research/v3/FOSP-A.md} (the fusion graph of
$\mathrm{SU}(2)$ is a one-dimensional path, so $\mathrm{TP}_2$ is automatic and carries no
information). The weight that the target measure actually induces is the frustrated one,
$(-1)^{2j}w_j$: although $n_{\text{negative,twisted}}=0$, the order does not preserve sign, and the
partial sums converge to $0.5034$ \LN. The audit's downgrade is accepted verbatim: what is proved
is that \emph{Lorentzian/M-convex has no referent in this framework} --- the support premise fails
and no such weight is constructed --- not that a mathematical object is impossible, which would
require a separate impossibility theorem not supplied by this wave \LN{} (framework). The
literature check adds that the raw material is missing as well: Lorentzian theory lives on
$\{0,1\}$ or M-convex supports, Wilson character coefficients change sign (Weingarten), and
stretched LR log-concavity is still unproved in general (\code{LIT-CHECK.md} Q2-4).

\textbf{What the table says as a whole.} Each of the five deaths is quantitative and each is
class-level: the number that kills it is a property of the mechanism, not of a script, so no future
tool built inside the class escapes it. In the charter's taxonomy all five are $(\gamma)$ results:
decisive refutations that tighten the map, not new species. That is the honest reading of a wave
with ``zero survivors'': five classes are now closed with receipts.

\subsection{The quantitative form of the data-mode theorem}

The meta-theorem of \code{research/unconditional/SPECIES-EXISTENCE.md} is a dichotomy. On a formal
class $\mathfrak C$ of candidate certificates, axiomatized by $F1$--$F4$, the computable subclass
$\mathfrak C^{\mathrm{comp}}=\mathfrak C_{\mathrm{uniform}}\cap\{F4\}$ is \emph{empty} ---
$F3\wedge F4=\varnothing$, ``$F3$'' being floor sensitivity and ``$F4$'' computability --- with the
evidence consisting of four data modes (finite moments, invariants, finite-$t$ analytic data,
quantitative extremal data), each of which is computable and floor-blind; while in the full class
$\mathfrak C_{\mathrm{uniform}}$ the existence of a certificate is \emph{equivalent} to the gap
itself. The first wave's contribution is to replace the emptiness statement by a number.

\begin{claimbox}{cT}
\claimhead{Definition}{(the data defect)} \emph{Fix a finite measure $\nu$ on $[0,\infty)$ with
$\nu([0,\infty))=1$ and a data map $D_N$ consisting of finitely many weak-$*$ continuous
functionals of $\nu$ --- canonically the first $N$ moments $M_0,\ldots,M_N$, and in the checks below
also the $L^1$ distance $D_8^{L1}$ on an eight-point grid with $M_1$ and $\chi$. Define}
\[
  d_N(\varepsilon):=\max\bigl\{\nu'([0,\varepsilon)):\ D_N(\nu')=D_N(\nu)\bigr\},
  \qquad
  d_N^{-}(\varepsilon):=\min\bigl\{\nu'([0,\varepsilon)):\ D_N(\nu')=D_N(\nu)\bigr\}.
\]
\emph{Thus $d_N(\varepsilon)$ is the largest mass that can hide below $\varepsilon$ without moving
the $N$-th order data at all, and $d_N^{-}$ the smallest.}
\end{claimbox}

\begin{claimbox}{cT}
\claimhead{Theorem}{(two-sided blindness, quantitative; scoped to the spectral class)} \emph{For
every finite $N$ the data fiber is non-trivial in both directions: it contains a measure with
strictly positive mass below $\varepsilon$ and a measure with zero mass below $\varepsilon$. Hence
neither ``the floor is $\ge\varepsilon$'' nor ``the floor is $<\varepsilon$'' is decidable from
$N$-th order data, at any finite $N$. The zero direction is elementary; the positive direction is
an explicit extremal construction, computed for all $N\le14$ \LE{} --- a finite computation, which is
evidence and not part of the theorem. The theorem is scoped to the
\emph{spectral class} --- finite lists of weak-$*$ continuous functionals of the K\"all\'en--Lehmann
measure. It says nothing about a data mode outside that class; constructing one is exactly the
charter's option $(\alpha)$, and this chapter does not claim to have found it.}
\end{claimbox}

The numbers are these (witness \code{w7/dm\_floor\_defect.py};
\code{ROUND1-VERDICTS.md} \S6). For the reference measure with $M_0=1$, floor $0.25$ and atom $0.4$
on a $301$-point grid, the positive direction holds at every $N\le14$:
\begin{align*}
  d_1(0.2)&=5.8219\times10^{-1}, & d_5(0.2)&=3.4479\times10^{-1},\\
  d_{10}(0.2)&=5.3758\times10^{-2}, & d_{14}(0.2)&=1.3634\times10^{-2},
\end{align*}
with $d_1(0.05)=5.601\times10^{-1}$ and $d_1(0.125)=5.717\times10^{-1}$ at the smaller windows; the
minimum is $0.0$ at every $N$ and every window \LT. A log-slope fit over $N=5,\ldots,14$ gives
$d_N\sim e^{-0.4595\,N}$ \LE{} (a fit on $N\le14$, not a theorem). Three corrections from the audit
are part of the statement and are kept
(\code{AUDIT-R1.md} \S C-5). First, the last three points turn \emph{up}
($d_{12}=1.0728\times10^{-2}$, $d_{13}=1.1192\times10^{-2}$, $d_{14}=1.3634\times10^{-2}$), so the
exponential is a fit and cannot be extrapolated; the honest content of the fit is only that the
defect stays small. Second, the assertion that ``buying accuracy by adding orders costs
$N=O(1/\varepsilon)$ orders'' is withdrawn, because it depended on the extrapolation. Third, the
\LT-grade part of the pair is the elementary zero direction; the positive direction is a computed
extremum. What survives is still the essential point: \textbf{a computable number now measures how
floor-blind a data mode is}, which is what turns the qualitative dichotomy into a quantitative one.

The second registered mode sharpens the same picture by failing faster. For the Laplace-sample mode
(samples $C(t)$ at $t\in[0,T]$ with $T=N/2$, step $0.5$) the mass that can hide below $0.2$ is
$4.879\times10^{-1}$ at $N=1$ and $2.183\times10^{-4}$ at $N=12$, again with minimum $0.0$
throughout \LE. More data of this kind is better at hiding, not worse at being blind.

The cleanest statement is the ray test, which is where the theorem becomes geometric. Put
\[
  \nu_h:=(1-h)\nu+h\,\delta_0,\qquad h\in[0,1].
  \label{eq:ray}
\]
Every weak-$*$ continuous datum moves \emph{linearly} in $h$ along the ray, while the floor
indicator $\mathbf 1_{G_{0.2}}$ --- the indicator of the event
$G_\varepsilon:=\{\nu:\inf\operatorname{supp}\nu\ge\varepsilon\}$ --- jumps only at the endpoint:
\begin{center}
\footnotesize
\begin{tabular}{@{}cccccc@{}}
\toprule
$h$ & floor mass & $D_8^{L1}(\nu_h,\nu)$ & $M_1$ & $\chi$ & $\mathbf 1_{G_{0.2}}$\\
\midrule
$10^{-1}$      & $0$ & $1.830\times10^{-1}$ & $1.603677$ & $1.653301$ & $0$\\
$10^{-2}$      & $0$ & $1.830\times10^{-2}$ & $1.764045$ & $1.818631$ & $0$\\
$10^{-3}$      & $0$ & $1.830\times10^{-3}$ & $1.780081$ & $1.835164$ & $0$\\
$10^{-4}$      & $0$ & $1.830\times10^{-4}$ & $1.781685$ & $1.836817$ & $0$\\
$10^{-6}$      & $0$ & $1.830\times10^{-6}$ & $1.781861$ & $1.836999$ & $0$\\
$10^{-9}$      & $0$ & $1.830\times10^{-9}$ & $1.781863$ & $1.837001$ & $0$\\
$0$            & $1$ & $0$                  & $1.7819$   & $1.8370$   & $1$\\
\bottomrule
\end{tabular}
\end{center}
That is, $D_8^{L1}=0.183\,h$ with three-digit constancy, $M_1\to1.7819$ and $\chi\to1.8370$ as
$h\downarrow0$, while $\mathbf 1_{G_{0.2}}$ is $0$ for every $h>0$ and $1$ at $h=0$ \LE{} for the
values, \LT{} for the structure: continuity of the data plus discontinuity of the indicator at the
ray endpoint. The structure is the same theorem as the weak-$*$ chapter's
no-continuous-forcing-functional result, instantiated on an explicit ray: any functional that
decides the floor must be discontinuous at $\nu$, and every finite list of continuous data is
therefore blind in \emph{both} directions --- it cannot certify the atom and cannot refute it. To
resolve the floor one needs the endpoint itself, i.e.\ an infinite-data limit, which is $H1$ again.

Two consequences follow, and they are the reason this subsection exists. (i) The meta-theorem's
emptiness is now \emph{quantified per data mode}: each mode has a defect profile, and the profiles
differ (moments $e^{-0.46N}$ by fit \LE; Laplace samples far faster), but every profile has minimum
zero, so no mode is rescued by adding order. (ii) The scope restriction is not cosmetic. The theorem
is about the spectral class, and a genuine fifth data mode --- computable \emph{and} floor-sensitive
--- is not excluded; it is the charter's $(\alpha)$ and it remains unbuilt. This is why the wave is
recorded as $(\gamma)$ with a quantitative body, not as a discovery.

\subsection{The verdict on $(NA)$, and the correction of its battlefield}

The named inequality $(NA)$ was recorded as the campaign's \emph{only variance-free face}:
\begin{equation}
  \operatorname{Cov}_\mu\bigl(G,\Phi'\mid0\bigr)\le0
  \qquad\Longleftrightarrow\qquad
  X:=-2\operatorname{Cov}_\mu\bigl(G,\Phi'\mid a_0\bigr)\ge0,
  \label{eq:naobj}
\end{equation}
with $V=S_W+\Phi$, $\Phi=-\log J$ the Haar Jacobian, $G=\partial_vV$ the force, and the
conditioning $\langle v,\xi\rangle=a_0$. If true it would pierce $H1$ without a variance estimate;
if it is not even statable, the last ``variance-free'' door is closed. The wave tested it in the
exactly solvable model
\begin{equation}
  S=\beta\Bigl(1-\cos\tfrac{\rho}{2}\cos\tfrac{\theta}{2}
  +\sin\tfrac{\rho}{2}\sin\tfrac{\theta}{2}\,u\Bigr),\qquad u=\frac{\xi_3}{\rho},\qquad
  J\propto\frac{\sin^2(\rho/2)}{\rho^2},\qquad \Phi=-\log J,
  \label{eq:namodel}
\end{equation}
i.e.\ one $\mathrm{SU}(2)$ plaquette with boundary holonomy $B=e^{i\theta\sigma_3/2}$, one free
link, longitudinal mode $v=e_3$, conditioning on $a_0$ along $e_3$ (model A of
\code{ROUND1-VERDICTS.md} \S7).

\textbf{Three methods, one number.} The covariance was computed by force moments and by finite
differences of the conditioned marginal; the two agree to $1.5\times10^{-13}$, and the
finite-difference curvature agrees with the analytic $G''$ to $1\times10^{-8}$ \LT/\LE. Quadrature
refinement is inert: node counts $n_s=1000,2000,4000,8000$ return exactly the same
$\operatorname{Cov}=-0.0158229063$, $X=+0.0316458126$. A three-method cross at the $10^{-13}$ level
in a model with one integration variable is as clean as this campaign's numerics get, so the
verdict below is not a numerical artefact.

\begin{claimbox}{cT}
\claimhead{Theorem}{($a_0=0$ is a zero set)} \emph{In model \eqref{eq:namodel}, for every $\theta$
and every $\beta$,}
\[
  X(a_0=0)=-2\operatorname{Cov}\bigl(G,\Phi'\mid a_0=0\bigr)\equiv0 .
\]
\emph{Proof.} $\Phi$ depends on $\xi$ only through $r=\lvert\xi\rvert$, so
$\partial_3\Phi=(\xi_3/r)\,[\,\cdot\,]$ is odd in $\xi_3$. The fibre $\xi_3=0$ is therefore a set on
which $\Phi'\equiv0$; hence $\operatorname{Var}(\Phi'\mid a_0=0)=0$ and the covariance vanishes for
every $G$. $\square$
\end{claimbox}

The theorem is the whole difficulty, and it cuts against the inequality rather than for it: the
\emph{literal} form of $(NA)$ at the natural conditioning point is $0\le0$ --- true, and empty. The
campaign's own audit had already recorded, in paraphrase, that the literal form is vacuous at the
natural conditioning point (\code{ROUND1-VERDICTS.md} \S7); what the wave adds is that the situation
reverses as soon as $a_0\ne0$.

\begin{claimbox}{cN}
\claimhead{Corollary}{($(NA)$ uniform in $a_0$ is false)} \emph{In model \eqref{eq:namodel} the
sign of $X$ depends on $a_0$. At $\theta=\pi$, $\beta=6$,}
\[
  X(a_0=-1)=-0.02307\ <\ 0\ <\ 0.10750=X(a_0=+1).
\]
\emph{Hence no symmetry, monotonicity, or reflection positivity argument can protect the sign of
the covariance, and the statement ``$(NA)$ holds uniformly in $a_0$'' is false.}
\end{claimbox}

The sign table, reproduced from \code{ROUND1-VERDICTS.md} \S7.3, is worth displaying because its
pattern is the result:
\begin{center}
\footnotesize
\begin{tabular}{@{}lccc@{}}
\toprule
 & $\theta=0$ & $\theta=\pi/2$ & $\theta=\pi$\\
\midrule
$a_0=0$  & $X\equiv-0.00000$ & $X\equiv-0.00000$ & $X\equiv-0.00000$\\
$a_0=-1$ & $+0.00182\ \ldots\ +0.00085$ & $-0.00119\ \ldots\ -0.00180$ & $-0.00430\ \ldots\ -0.02307$\\
$a_0=+1$ & $+0.00182\ \ldots\ +0.00085$ & $+0.00429\ \ldots\ +0.01025$ & $+0.00493\ \ldots\ +0.10750$\\
\bottomrule
\end{tabular}
\end{center}
Each entry is a range over $\beta=0.5,\ldots,6$. Three readings: at $a_0=0$ the value is zero for
all $\theta$ and $\beta$ (\LT, the theorem above); at $a_0=-1$, $\theta\in\{\pi/2,\pi\}$ it is
\emph{negative for every $\beta$} in the scan, i.e.\ $X\ge0$ fails; and the same $\theta$ with
$a_0=+1$ gives a positive value that grows with $\beta$. The magnitudes are $10^{-2}$ to $10^{-1}$,
ten to twelve orders of magnitude above the quadrature agreement
($10^{-2}/1.5\times10^{-13}\approx7\times10^{10}$), so the sign flip is not marginal \LE.
Consequently $(NA)$ is available in exactly one of two forms, and both are dead: pointwise at
$a_0=0$ it is vacuous, and uniform in $a_0$ it is false \LN.

\textbf{The missing constant, independently reproduced.} The same model contains the ultraviolet
chapter's counterexample, and the wave reproduced it without using the earlier code. At
$\theta=\pi$, $v=e_3$, $a_0=0$,
\begin{equation}
  G''=A-B\beta^2,\qquad A=0.20511530,\quad B=0.01105929,\quad \beta_*=4.30661,
  \label{eq:nacounter}
\end{equation}
so $\beta=2.5$ gives $+0.13599$ and $\beta=6$ gives $-0.19302$, and the curvature changes sign at
$\beta_*=4.30661$; all five digits agree with \code{research/v3/V3-COUNTEREXAMPLE-SCOPE.md} \LE. The mechanism
is also reproduced, and it is the mechanism that matters for this chapter: conditioning at $a_0=0$
squeezes $E[H\mid a_0=0]$ to \emph{exactly} zero, so the curvature reduces to
$G''=+W_{\mathrm{geo}}-\operatorname{Var}$, and the variance grows with $\beta$ ($0.0691$ at
$\beta=2.5$, $0.3981$ at $\beta=6$) until the negative term wins. A cross-check by an independent
route is worth more than a new conjecture; this one also explains \emph{why} the failure is at
$\beta>\beta_*$ and not at small coupling.

\textbf{The structural observation: the battlefield is $\mathrm{SU}(N\ge3)$.} The theorem above is
not a fact about Yang--Mills; it is a fact about one-plaquette $\mathrm{SU}(2)$. For $\mathrm{SU}(2)$
the Jacobian of the Haar measure in exponential coordinates is a function of $r=\lvert\xi\rvert$
alone, because the Weyl group is $\{\pm1\}$ acting by $\xi\mapsto-\xi$, and a Weyl-invariant
function of a three-vector in the fundamental representation can only depend on its length. Radial
symmetry is exactly what makes $\partial_3\Phi$ odd and the $a_0=0$ fibre a zero set. For
$\mathrm{SU}(N\ge3)$ the Jacobian is invariant only under the Weyl group of the Cartan torus --- a
finite reflection group acting on the $N-1$ Cartan coordinates, not a rotation of the whole Lie
algebra --- so $J$ is \emph{Weyl-symmetric but not radial}, the fibre $\langle v,\xi\rangle=0$ is not
a vanishing set of $\partial_v\Phi$, and the orbital integral does not degenerate. Concretely, the
derivative that controls the covariance is
\begin{equation}
  \partial_v I_3=3\operatorname{tr}\bigl(T_v\,(\xi_\perp\!\cdot T)^2\bigr)\neq0 ,
  \label{eq:nasun}
\end{equation}
which vanishes identically in the $\mathrm{SU}(2)$ model only because there is a single
$\mathrm{SU}(2)$-invariant quadratic in $\xi$ (the length). Hence:
\begin{qbox}
\textbf{The structural correction.} The identity $X\equiv0$ at $a_0=0$ is an accident of
$\mathrm{SU}(2)$ radiality. The true content of $(NA)$ --- if any --- lives on
$\mathrm{SU}(N\ge3)$, where $J$ is only Weyl-symmetric and $\partial_vI_3\neq0$. The corpus tested
the inequality in precisely the one model in which its natural form is trivial, and the only
$\mathrm{SU}(N\ge3)$ test performed since is the one-plaquette model reported at the end of this
chapter \LN{} for ``$(NA)$ is the variance-free face''; the $4$D $\partial_vI_3$ source is open and is
\emph{not} claimed closed here.
\end{qbox}

\textbf{Why the class is closed twice over.} The literature check (\code{LIT-CHECK.md} Q1--Q1$'$)
supplies the second closing. (i) No published negative-association or anti-FKG theorem covers a
Haar\,$\times$\,Wilson measure: the whole machinery (Pemantle; Borcea--Br\"and\'en--Liggett;
Joag-Dev--Proschan; Harg\'e) lives on lattice orders, polynomials, log-concave densities, or
Gaussian structure, and a compact group with a periodic log-density has none of the four \LT{} (a
statement about the literature, negative in form). (ii) The one general sign theorem that
\emph{does} exist points the other way: the Stein/score identity gives
$\operatorname{Cov}(\partial_iS,\partial_jS)=E[\operatorname{Hess}S]_{ij}$, so there is no
universal sign law, and the implication ``$\operatorname{Hess}S\succeq\rho\Rightarrow$ gap
$\ge\rho$'' is Bakry--\'Emery $\Gamma_2$ --- which fails here \emph{by construction}, because the
Savvidy Hessian $-D^2+2\operatorname{ad}F_B$ is indefinite (the negative mode is transverse
magnetic). Therefore $(NA)$, read correctly, is a \emph{Hessian-sign} statement, i.e.\ a statement
that the variance term wins in $G''=E[H]-\operatorname{Var}+W_{\mathrm{full}}$ --- and that is $H1$
in another notation \LC. The label on the class is \LN: the face is not variance-free at all; it
was $H1$ with the variance moved to the other side of the inequality.

\textbf{What the wave achieved.} One new no-go of the sharpest kind (a named inequality is either
vacuous or false, decided by a two-line oddness argument plus a sign table), one independent
five-digit reproduction of the ultraviolet counterexample, and one structural relocation of the
question: from ``$(NA)$ on $\mathrm{SU}(N)$'' to ``$\partial_vI_3\neq0$ on $\mathrm{SU}(N\ge3)$''.
The last is the only place where a variance-free argument could still live, and the campaign's
model for it is the one-plaquette $\mathrm{SU}(3)$ test recorded at the end of this chapter.

\subsection{Two load-bearing constants corrected}

This subsection is the chapter's erratum interface. The main text is edited elsewhere; what is
recorded here is the correction, its evidence, and its price. Both constants are load-bearing in the
conditional theorem, which is why a wrong value is worse than a gap: it propagates silently into
every downstream consequence.

\textbf{(A) The ratio constant: $C=2$, not $C\approx1.5$.} The conditional chain uses a window on
the second spectral gap,
\begin{equation}
  \rho^2:=\frac{\mu_2}{\mu_1}\ \le\ C^2 ,
  \qquad \mu_0=0<\mu_1\le\mu_2\le\cdots ,
  \label{eq:window}
\end{equation}
where $\mu_n$ are the eigenvalues of the $0^{++}$-projected coarse transfer generator and
$\mu_1,\mu_2$ are \emph{distinct} values (the K\"all\'en--Lehmann measure merges coincident masses
into one atom). The corpus's \code{research/unconditional/GAP-IN11.md} (\code{IN-11}) used
$C\approx1.5$, i.e.\ $\rho^2\le2.25$. That value is
violated by five of the exact anchors below --- all but Mehler $2$ and 2D $\mathrm{SU}(3)$ $9/4$,
which sits exactly on the bound (\code{RB-RATIO-BOUND.md} \S RB-8):

\begin{center}
\footnotesize
\begin{tabular}{@{}lccc@{}}
\toprule
\textbf{Exact anchor} & $\rho^2$ & $C=1.5$ ($C^2=2.25$) & $C=2$ ($C^2=4$)\\
\midrule
Mehler / Ornstein--Uhlenbeck heat kernel & $2$ & pass & pass\\
one plaquette $\mathrm{SU}(2)$, $\beta=0$ ($\mu_1=3$, $\mu_2=8$ exactly) & $8/3=2.6667$ & \textbf{fail} & pass\\
2D Yang--Mills $\mathrm{SU}(2)$ & $8/3=2.6667$ & \textbf{fail} & pass\\
2D Yang--Mills $\mathrm{SU}(3)$ & $9/4=2.25$ & pass (boundary) & pass\\
path graph, $N=2$ & $3$ & \textbf{fail} & pass\\
box / 1D rotor / 1D Dirichlet string & $4$ & \textbf{fail} & pass (\textbf{saturated})\\
2D $U(1)$ & $4$ & \textbf{fail} & pass (\textbf{saturated})\\
\bottomrule
\end{tabular}
\end{center}

Five exact anchors violate $C=1.5$; two of them saturate $C=2$ exactly \LT. The corpus's value is not
arbitrary --- it is precisely the 2D $\mathrm{SU}(3)$ Casimir ratio $9/4$ --- which is presumably
how it was chosen, and which is also why it cannot be a class-independent constant: the family moves away
from $9/4$ in both directions ($8/3$ at $\mathrm{SU}(2)$, $4$ at $U(1)$ and at the
Dirichlet string). This is \LT{} for the exact rows and \LE{} for the finite-difference
confirmation (\code{RB-RATIO-BOUND.md} \S RB-8, PART E--F).

The 1D anchor also tells us \emph{which} convention is load-bearing. For $-u''+qu=\lambda u$ on
$[0,1]$ with Dirichlet data the eigenvalues are $\mu_k=k^2\pi^2+q$; writing the constant potential as
$q=-c$ (with $0\le c<\pi^2$) and as $q=+c$, the two closed forms are
\begin{equation}
  \rho^2_{-}(c)=\frac{4\pi^2-c}{\pi^2-c}\quad(q=-c),
  \qquad
  \rho^2_{+}(c)=\frac{4\pi^2+c}{\pi^2+c}\ \le 4\iff c\ge0\quad(q=+c,\ \text{sharp}) ,
  \label{eq:dirichletratio}
\end{equation}
and both signs must be kept because they behave oppositely: $\rho^2_{+}$ decreases in $c$ and
saturates $4$ exactly at $q\equiv0$, while $\rho^2_{-}$ diverges as $c\nearrow\pi^2$. So $C^2=4$ is
exactly the Dirichlet-string saturation value \LT. The bound is sharp and the
mechanism transparent: the sharp half $q\ge0$ means the effective potential's bottom is not pushed
below the free value, i.e.\ $e_1\ge\pi^2/L_s^2$ --- a \emph{floor}, not a degree condition. The
finite-difference check at $N=4000$ confirms the $q=-c$ branch of \eqref{eq:dirichletratio} ($c=9.7$:
exact $175.5757$, FD $175.5762$), and a $215$-shape battery with $q\ge0$ gives $\sup\rho^2=3.999999$
(attained at $q\equiv0$) and $\inf\rho^2=1.00000$ (a symmetric barrier, i.e.\ an exactly degenerate
doublet) \LE.

The floor, however, is one-sided, and one-sided is not enough. \code{research/v3/species/DT-DIRICHLET.md}
records two exact witnesses that $e_1\ge\pi^2/L_s^2$ on its own supplies \emph{no} ratio bound.
\emph{Witness B} is local and sign-changing with the floor already satisfied:
$e_1=18.0963\ge\pi^2$ yet $\rho^2=27.1846$, so a floor at the bottom does not stop the second level
from running away. \emph{Witness A} is a rank-one positive bump that moves only the second level,
$\rho^2=4+M/\pi^2=4.5,\,9.0$, and lifting the first $r$ levels together gives
$\rho^2=4,9,16,24,\dots$, unbounded. The legal form of the constant is therefore a \emph{two-sided}
comparison, $c_-L_0\le L\le c_+L_0\Rightarrow C=2\sqrt{c_+/c_-}$, and the window is a floor
proposition only once the trace input is attached: with $t:=T_2(V)$, $\rho^2\le C^2$ holds if and
only if $T_2(V)\le(1+C^2)e_1$, which at $C=2$ reads $T_2(V)\le5e_1$ \LT{} (the algebra) \LE{} (the
witnesses). The general
statement ``$q\ge0\Rightarrow\rho^2\le4$'' has a numerical battery and a proof route
(ground-state transform plus Pr\"ufer comparison) but no proof, so it is \LC, not \LT.

\begin{claimbox}{cC}
\claimhead{Proposition}{(the best constant inside the Dirichlet class)} \emph{The largest constant
that survives every exact anchor \textbf{inside the Dirichlet-type class} (bottom-normalised
effective potential, single well, non-degenerate bottom) is $C=2$, i.e.\ $\rho^2\le4$, with equality
attained by the 1D Dirichlet string and by 2D $U(1)$; \textbf{outside that class it fails} --- the
Neumann--Dirichlet box gives $\rho^2=9$ exactly and the star $K_{1,m}$ gives $\rho^2=m+1$,
unbounded \LC. In the window form \eqref{eq:window} the price is a weaker floor consequence:}
\[
  \mu_1\ \ge\ \frac{\Delta}{C^2-1}:\qquad
  \Delta/1.25\ \longrightarrow\ \Delta/3 .
\]
\emph{The window theorem itself is unchanged; only its constant and the floor it yields move.}
\end{claimbox}

\textbf{Why the correction is forced by the physical spectrum, not only by solvable models.} The
decisive argument is not a toy: the physical ratio lies in the band $1.54$--$1.56$ reported in
\code{research/unconditional/RATIO-RHO.md} \S5 (the older $1.6$--$1.8$ band of
\code{research/v3/Z1-HOLO.md} gives larger values still), so
\begin{equation}
  \rho^2_{\mathrm{phys}}\in[1.54^2,\,1.56^2]=[2.3716,\,2.4336],
  \label{eq:physwindow}
\end{equation}
and \textbf{every point of the band already exceeds $2.25$}, which is why the exclusion of $C=1.5$
is robust to which end is taken. The corpus's window was therefore violated by the theory it was
supposed to constrain, before any model of it was written down. With $C=2$ the whole band sits
comfortably inside the bound, which
is what a class-sharp constant does: exclude the gapless configurations of the class while admitting
the physical one \LE{} for the lattice ratio, \LT{} for the arithmetic. The trade is honest: a
weaker consequence from a true constant beats a stronger consequence from a false one.

\textbf{What $C=2$ is, and what it is not.} $C=2$ \emph{is} the sharp value inside the Dirichlet-type
tower class (bottom-normalised effective potential, single well, non-degenerate bottom), saturated by
the 1D Dirichlet string; it is \textbf{not} a structural constant: Neumann/mixed boundary conditions
give $\rho^2=9$ (box, exact) and $\rho^2=m+1$ (star $K_{1,m}$, unbounded), and the same battery
records $\rho^2=6.6928$ for the Neumann--Dirichlet box at $q=1$ and $5.2689$ for the 2D disc's
$s$-wave sector (\code{research/v3/species/w7/kt\_a\_exact.log} A8). Hence $\rho^2\le C^2$ and the
missing step --- the Dirichlet normalisation of the confining coordinate --- are \textbf{the same
condition, not independent inputs}. The floor is $\Delta/(\rho^2-1)$: $\Delta/3$ inside that class,
degrading outside --- which is precisely the missing step, not an independent input.

\textbf{(B) The strong-coupling radius: $1/48$, not $1/24$.} The second correction is already
recorded in the errata appendix and is cross-referenced here rather than re-argued
(Appendix~\ref{app:erratum-radius}). Shen--Zhu--Zhu prove the mass gap for
$\lvert\beta\rvert<1/(16(d-1))$, which in $d=4$ reads $\lvert\beta\rvert<1/48$ and \emph{does not
involve $N$}; written in the 't Hooft coupling it is $\lambda_{\mathrm{tH}}=\beta N<N/48$, which is
a different and $N$-dependent statement. The inequality $\lambda_{\mathrm{tH}}<1/24$ occurs nowhere
in the source. The provenance is diagnosable --- $1/24=2\times1/48$, i.e.\ a factor-two convention
slip (or a doubled rumour of the CNS improvement) --- and no conclusion depends on it: the physical
coupling $\beta\approx2.5$ is $120\times$ the corrected $\beta$-radius, and in 't Hooft units at
$\lambda_{\mathrm{tH}}=1$ the corrected radius is $48/N$, i.e.\ $24\times$ ($N=2$) or $16\times$
($N=3$) away. The rigorous frontier remains two orders of magnitude from the physical point
\LT{} (source check) plus the unit arithmetic.

\textbf{What the two corrections have in common.} Both constants were doing work in the conditional
theorem, both were plausible, and both were wrong in the direction that \emph{flatters} the
argument: $1.5$ made the window look tighter than the spectra allow, and $1/24$ made the rigorous
frontier look closer than it is. Neither error touched the verdict --- the gap is still open either
way --- but both would have corrupted any future use of the chain. Recording them is part of the
ledger, and it is the reason this chapter exists as a separate section rather than as a footnote in
the chapters that consume the constants.

\subsection{(TCI) is under-specified, not merely unproved}

The third named inequality, at the level of its own definition, is a template with holes. As printed
it reads $\langle\varphi,Q\varphi\rangle+\langle\varphi,\mathbb K\varphi\rangle-cF\ge\mu>0$ on
$\mathcal H_-$ and defines \emph{none} of $Q$, $\mathbb K$, $cF$, $\mathcal H_-$
(\code{TCI-COMPENSATION.md} \S1.1). Concretely: $cF$ is fixed only outside the paper, as
$cF:=g\lvert B\rvert$ (\code{research/v3/SM5-TCI-HEAD.md}:74,91); the letter $Q$ carries two inequivalent
meanings, the interaction second variation (implied by ``$Q=0$ on the free massless field'') and the
Savvidy operator $Q_{\mathrm{Sav}}=-D^2+2\operatorname{ad}F_B$ (\code{research/v3/SM5-TCI-HEAD.md}:72); and
$\mathbb K$ is a promissory note --- ``the lift of the block scheme'', with no formula anywhere.

The only self-consistent completion is $Q=3\lambda_4\varphi_0^2=3\,cF$ together with the
seven-element enumeration $\mathbb K\in\{0,\,-BC^{-1}B^{*},\,\mathrm{Hess}\,\Phi,\,(d-3)cF,\,
c(-\nabla^2)^p,\,\Gamma_2,\,0_{\mathrm{top}}\}$ \LC{} (a choice: the source labels
$Q=3\lambda_4\varphi_0^2=3\,cF$ as $[MINE]$ and the enumeration as its one non-obvious option
$(\mathrm{D}3)$), under which the inequality becomes
$(\mathrm{TCI})\iff\lambda_{\min}(A_{\mathbb K}\mid\mathcal H_-)\ge cF+\mu$ \LT. Two consequences
follow, and the first is a strict non-equivalence that the corpus had conflated:
$(\mathrm{TCI})\not\Rightarrow\Sigma\succeq\mu$, because $\Sigma\preceq A_{\mathbb K}$ always, so
the Schur form is strictly stronger and the gap between them is exactly $\lVert B\rVert^2/\gamma$.

Compensation is nevertheless not elimination, and this is the species' real content \LT:
$\Sigma\preceq A$ holds always, so elimination requires $\gamma>0$ and can only \emph{lower} the
in-block form, whereas compensation is $\gamma$-free and only raises it. On the $4$D $L=1$ scan the
in-block curvature is $+0.1667\to+0.7280$ and monotone in the field for $c=0\to1.5$, while $\Sigma$
decays to $+0.0109$ at $c=0.52$ and the Schur form ceases to exist at $c_\gamma\approx0.52$; and
every Bogoliubov/symplectic lift is sign-negative by construction, since $\omega^2\to\omega^2-4c^2$
(ground energy $1.00000\to0.97980\to0.91652\to0.60000$ for $c=0,0.1,0.2,0.4$). What kills the
species is not the machine but the source. On a $2\times2$ block, $\mathrm{SU}(2)$, $\beta=2.5$,
$6\times10^4$ Haar samples, $\dim\mathcal H=12$: $\mathrm{Hess}\,\Phi\succeq(N/12)\mathbb 1$
\emph{exactly} ($0.166680$ against $0.166667$); $E[n_-]=1.70$ with $P(n_->0)=0.851$; the
infimum over $\mathcal H_-$ has median $-0.1796$ ($q_{05}=-0.6745$); and $P(\mu_3\ge0)=0.070$,
short by $2.72\times$, with $\beta_c\approx0.76$ \LE. The compensated well exists but sits at
$a_*\approx9$--$14$, outside any small-block amplitude, so no small-volume witness can see it \LE.

\begin{claimbox}{cN}
\claimhead{Verdict}{(TCI) is H1 in another coordinate} \emph{Because $V=S_W+\Phi$ has exactly two
additive pieces, any second-variation term is a piece of those two: the bounded Haar/Jacobian floor
$N/12$ (measured $2.72\times$ too small at the head) or the condensate $(d-3)cF$, which is the
target itself. There is no third positive order-$0$ source, so the only non-vacuous form of
$(\mathrm{TCI})$ is $(d-3)cF\ge\mu$, i.e.\ $\lambda_4\langle\lvert a_0\rvert^2\rangle\ge\mu$ ---
the uniform positivity of the gluon condensate.}
\end{claimbox}

The label is \LN{} for the class: $(TCI)$ was never merely unproved, it was unspecified, and its
one well-defined reading is $H1$ after a change of variables that moves the condensate from the
right-hand side to the left. A separate formulation defect is recorded with it: $\mathcal H_-$ is
defined \emph{by} the Hessian, so on the free massless field $\mathcal H_-=\{0\}$ and the literal
$\forall\varphi\in\mathcal H_-$ form holds \emph{vacuously} rather than failing the acid test ---
which is the opposite of a pass \LT.

\subsection{Two bridges: from $W$ to $I$, and from $I$ to a computable threshold}

The second wave accepted the charter's direction and stopped hunting inside the variance/KL category.
It attacked the campaign's own two remaining targets: the \emph{alignment} predicate
$W:\,w_1=f_1^2/M_0\ge w_0$, and the \emph{ratio} predicate behind the window,
$\rho^2\le C^2$. The logical geography is fixed (\code{research/v3/U6-BINDING-LOGIC.md},
\code{research/v3/T6-BINDING-ADVERSARIAL.md}): $W\Rightarrow A$, $I\Rightarrow B_{\mathrm{spec}}\Rightarrow A$,
$W\perp I$, $W\wedge H1\Rightarrow G$, and $I\wedge(\rho^2\le C^2)\Rightarrow G$, where $A$ is
$f_1^2>0$, $I$ is the isolation $\Lambda^\star=\inf\operatorname{supp}(\nu-f_1^2\delta_{m_1^2})>m_1^2$,
and $G$ is the floor $m_1^2>0$. Both attacks return the same shape: a true identity that points the
wrong way, one true lower bound that is sharp and nevertheless contains no Yang--Mills information,
and a computable half whose other half is a comparison operator --- i.e.\ confinement dynamics.

\textbf{Bridge one: an exact identity that goes the wrong way \LT.} Let $T=e^{-aH}$ be the
reflection-positive transfer matrix, $\Omega$ the ground state, $o$ a bounded self-adjoint source
with $\langle o\rangle_\mu=1$, and $w_n=|\langle\psi_n|o|\Omega\rangle|^2/\operatorname{Var}_\mu(o)$,
$s_n=E_n-E_0$, so that $\sum_{n\ge1}w_n=1$. Then
\begin{equation}
  \sum_{n\ge1}(E_n-E_0)\bigl\lvert\langle n|o|\Omega\rangle\bigr\rvert^2
  =\tfrac12\langle\Omega|[o,[H,o]]|\Omega\rangle
  =\bigl\lVert (H-E_0)^{1/2}(o-\langle o\rangle)\Omega\bigr\rVert^2\ \ge\ 0 ,
  \label{eq:dcom}
\end{equation}
and in normalized form
\begin{equation}
  \bar E:=\sum_{n\ge1}s_nw_n=\frac{1}{2}\frac{D_o}{\operatorname{Var}_\mu(o)},
  \qquad D_o:=\langle\Omega|[o,[H,o]]|\Omega\rangle ,
  \qquad \bar E=-\frac{d}{dt}\log C(t)\Big|_{t=0}.
  \label{eq:barE}
\end{equation}
The proof is two lines --- $[o,[H,o]]=2oHo-o^2H-Ho^2$, take the expectation in $\Omega$, use
$H\Omega=E_0\Omega$ and the spectral expansion --- and the identity is exact, not an estimate \LT.
It was checked on the double well to a relative residual $4.3\times10^{-5}$ (truncation in $k$) with
completeness $\sum_n|\langle n|x|0\rangle|^2=\operatorname{Var}x$ to $1.3\times10^{-15}$ \LE.

The identity is the campaign's cleanest negative result, because its \emph{direction} is fixed by
its being a first moment. Since $s_n\ge s_1$ we get $w_1\le\bar E/s_1$ and
$\sum_{n\ge2}w_n\le\bar E/s_2$: both are \emph{upper} bounds, and the first needs the gap it was
supposed to help prove. Conversely, with $(M_0,\bar E)$ held fixed, $w_1$ sweeps the whole interval
$[0,\bar E/s_1]$: for any $\varepsilon\in[0,\bar E/s_1]$ the two-point measure
$\nu=\varepsilon\delta_{s_1}+(1-\varepsilon)\delta_{s^*}$ with $s^*=(\bar E-\varepsilon s_1)/(1-\varepsilon)$
has the same data and $w_1=\varepsilon\to0$ \LT/\LE. Hence \textbf{moments are an upper-bound
machine, and $w_1\ge w_0$ cannot be proved from any finite set of them} \LN. This is the same
statement as the data-mode theorem of the previous subsection, seen from the source side: the moment
defect of the data-mode theorem above is the quantitative form of exactly this witness.

\textbf{Bridge one, continued: the only lower bound, and what it costs \LT.} One inequality does go
downwards, the residue--level-gap lemma. If $\nu\ge0$, $\nu([0,\infty))=1$, $s_1=\inf\operatorname{supp}\nu$,
and $\nu$ has no mass in $(s_1,\Lambda)$, then
\begin{equation}
  w_1\ \ge\ \frac{\Lambda-\bar E}{\Lambda-s_1},
  \label{eq:l2}
\end{equation}
with equality for the two-point measure $\varepsilon\delta_{s_1}+(1-\varepsilon)\delta_\Lambda$ ---
so the bound is sharp, and being sharp is precisely why it cannot be a mechanism. At $\Lambda=s_2$
it is the truncation of an \emph{identity}, and the identity is the sharper object
(\code{research/v3/species/QI-BRACKET.md} QI-1):
\begin{equation}
  w_1\ =\ \frac{s_2-\bar E}{s_2-s_1}\ +\ \frac{\Delta}{s_2-s_1},
  \qquad
  \Delta:=\int_{[s_2,\infty)}(s-s_2)\,d\nu\ \ge\ 0 ,
\end{equation}
and \textbf{the inequality form \eqref{eq:l2} is the $\Delta\ge0$ truncation of this identity}. The
slack $\Delta/(s_2-s_1)$ is the UV-sensitive excess first moment of the tail above $s_2$
(not the transfer gap $\Delta$ of the tropical subsection); it is exactly the piece that a
level-counting certificate cannot see, which is why the bridge is a few-level certificate and not a
mechanism. Substituting
\eqref{eq:barE} turns the wall into a single comparison between one computable number and one
support point:
\begin{equation}
  w_1\ge w_0
  \qquad\Longleftarrow\qquad
  s_2\ \ge\ \frac{\bar E-w_0s_1}{1-w_0},
  \label{eq:bridge}
\end{equation}
where $s_2$ is the \emph{second support point} of the measure (not the second energy level), and
$\bar E$ is computable. The qualifier ``second support point'' is not pedantry: taking
$\Lambda=1.3\,s_2$, i.e.\ past the second support point, makes \eqref{eq:l2} over-report in all
seven couplings tested ($0.9935>0.9911$, $0.9868>0.9822$, $0.9664>0.9562$, $0.9560>0.9437$,
$0.9596>0.9488$, $0.9712>0.9635$, $0.9778>0.9716$) \LE. So $1.3\,s_2$ cannot be used as a proxy for
$s_2$: a certificate built on it is a false certificate. The degenerate limit is the free field:
there the measure is a pure cut, no second support point exists, the denominator of
\eqref{eq:bridge} loses its second node, and the inequality is not weak but undefined --- which is
how the route passes the free-field acid test without saying anything about Yang--Mills.

\textbf{Bridge two: a computable rank-two trace \LT.} The ratio coordinate has an exact
variational identity after Ky Fan. With $L\ge0$ self-adjoint, discrete spectrum, and
$\ker L=\operatorname{span}\{1\}$ simple,
\begin{equation}
  \begin{gathered}
    \mu_1+\mu_2=\min\bigl\{\operatorname{tr}(P_VLP_V):\ \dim V=2,\ V\perp1\bigr\},\\
    T_2(V):=\operatorname{tr}(G^{-1}M),\qquad
    G_{ij}=\langle\varphi_i,\varphi_j\rangle,\qquad
    M_{ij}=\langle\varphi_i,L\varphi_j\rangle ,
  \end{gathered}
  \label{eq:kyfan}
\end{equation}
so that for \emph{every} test pair
\begin{equation}
  \rho^2=\frac{\mu_2}{\mu_1}\ \le\ \frac{T_2(V)}{\mu_1}-1 .
  \label{eq:ratio}
\end{equation}
Three properties make \eqref{eq:ratio} the useful half. It is independent of $\psi_1$ (the two
dimensionality absorbs the $\perp\psi_1$ constraint automatically); it contains no small parameter
and no $\beta$ expansion, so $t:=T_2(V)$ is an $(L,a)$-uniform pair of numbers once $V$ is fixed;
and its direction is right --- an upper bound on a minimum is obtained by \emph{exhibiting} a test
space, which is constructive. Numerically, on a path graph with $N=40$ ($\mu_1=0.023439$,
$\mu_2=0.052609$), two thousand random test spaces give $T_2(V)-(\mu_1+\mu_2)\in[2.529,5.420]>0$
with no violation, and $V=\operatorname{span}(\psi_1,\psi_2)$ returns
$T_2=0.076048305=\mu_1+\mu_2$ to $10^{-9}$ \LE. The trace is also \emph{additive} over commuting
generators, $T_2(V)=T_2^{(1)}(V)+T_2^{(2)}(V)$, which is why it survives a block-spin step and is
in practice computable: no other object in the registry is an additive functional of the generator.

\textbf{Bridge two, continued: the equivalence, and the no-go \LN.} Writing
$\eta:=\mu_1/(t/(1+C^2))$, the window is exactly
\begin{equation}
  \rho^2\le C^2
  \iff \mu_1\ \ge\ \frac{t}{1+C^2}
  \iff \eta\ge1 ,
  \label{eq:eta}
\end{equation}
a \emph{floor} proposition with a computable, explicitly dimensioned threshold \LT. This is the
sense in which $\rho^2\le C^2$ is not a fourth independent predicate: it inherits $H1$ \LN{} for the
named class ``ratio-coordinate claims''. The reason is structural and was proved as a direction
theorem: every min--max, Ritz, Ky Fan, interlacing, or second-cone device returns upper bounds on
eigenvalues and $\sup_V\nu_k(V)=+\infty$, so test spaces can never lower-bound $\mu_1$; the only
source of a lower bound is a \emph{formal comparison} $L\ge L_0$ with $e_1(L_0)=\lambda$, and that
is a statement that the effective potential has no lower well. The ratio coordinate therefore
decomposes into ``construction $+$ comparison'', and only the construction half was ever missing
from the toolbox.

\textbf{The acid test both bridges must pass, and do not.} The double well
($H=-\tfrac12d^2/dx^2+(g/2)(x^2-1)^2$, source $o=x$, so $[x,[H,x]]=1$ and
$\bar E=1/(2\operatorname{Var}x)$ exactly for every $g$) is the model in which alignment is true and
the gap is not:

\begin{center}
\footnotesize
\begin{tabular}{@{}ccccccc@{}}
\toprule
$g$ & $s_1=E_1-E_0$ & $s_2$ & $w_1$ & $\bar E_{\mathrm{spec}}$ & L2 bound & $\chi=\sum w_n/s_n$\\
\midrule
$0.5$   & $7.250\times10^{-1}$ & $3.4438$  & $0.991131$ & $0.749266$ & $0.9911$ & $9.14\times10^{-1}$\\
$2.0$   & $7.918\times10^{-1}$ & $4.7958$  & $0.982232$ & $0.863557$ & $0.9821$ & $7.20\times10^{-1}$\\
$8.0$   & $4.624\times10^{-1}$ & $6.3232$  & $0.956187$ & $0.723110$ & $0.9555$ & $1.435$\\
$16.0$  & $1.914\times10^{-1}$ & $7.5450$  & $0.943689$ & $0.614551$ & $0.9425$ & $4.018$\\
$32.0$  & $3.887\times10^{-2}$ & $9.9225$  & $0.948802$ & $0.558846$ & $0.9474$ & $2.185\times10^{1}$\\
$64.0$  & $3.033\times10^{-3}$ & $14.2506$ & $0.963456$ & $0.536284$ & $0.9626$ & $2.962\times10^{2}$\\
$100.0$ & $3.015\times10^{-4}$ & $18.2581$ & $0.971625$ & $0.527945$ & $0.9711$ & $3.052\times10^{3}$\\
\bottomrule
\end{tabular}
\end{center}

Three readings, each of which is a verdict. (i) $w_1$ stays in $0.94$--$0.99$ while $s_1$ collapses
by $3.4$ orders of magnitude and $\chi$ grows from $0.91$ to $3.1\times10^3$: alignment is
\emph{blind} to the floor, so a lower bound on $w_1$ that holds here carries no gap information
\LE. (ii) The L2 bound is saturated, with slack $w_1-\mathrm{L2}\in[4.7\times10^{-5},
1.4\times10^{-3}]$ \LE: the inequality is not wrong, it simply fires where it is not needed, which
is the precise definition of a mechanism that is not a mechanism. (iii) A different double well
kills any \emph{upper-bound-data} machine for the ratio: the SUSY form of
\code{RB-RATIO-BOUND.md} \S RB-7 ($H=-\partial_x^2+W$, $W=V'^2/4-V''/2$, $V=(g/2)(x^2-1)^2$,
FD $N=6000$) --- \emph{not} the $H=-\tfrac12d^2/dx^2+(g/2)(x^2-1)^2$ of the table above whose
$s_1=0.7250$ belongs to another operator. For it, as $g$ runs $0.5\to12$ the ratio grows
$\rho^2=4.48\to807.36$ (with $\log\rho^2\approx0.558g$, i.e.\ exponentially) while
$T_2=\mu_1+\mu_2$ grows only $4.34\to20.05$ --- a factor $180$ against a factor $4.6$ \LE.
Hence $\rho^2$ is not a continuous functional of upper-bound data; its explosion is entirely the
Kramers collapse of $\mu_1$.

\textbf{The numerical caution that keeps the bridge honest.} Feeding real lattice data into
\eqref{eq:l2} currently produces a \emph{false} certificate rather than a weak one. On the existing
$\beta=2.5$, $L=4,6,8$, $N=2000$ correlators the $0^{++}$ channel is contact/zero-mode dominated
($|C(t)/C(0)-1|\le4.7\times10^{-3}$, non-monotone), so $\bar E_{\mathrm{bare}}\approx+0.005/
-0.002/-0.001$, and \eqref{eq:bridge} then ``proves'' $w_1\gtrsim0.99$ --- a certificate made of
normalization noise. The diagnostic is exact: replacing $C(t)$ by a completely flat constant and
refitting with the same form and the same $m_1=0.03$ reproduces $A/C(0)=0.5308/0.5467/0.5629$
against the measured $0.5288/0.5475/0.5633$, so the quoted amplitude contains no $f_1^2$
information at all \LE. The bare total normalization also grows super-extensively,
$M_0/V_s=206\to1170\to3894$ for $L=4,6,8$, confirming that the ratio built on it is an ultraviolet
artefact. The bridge is therefore correct as mathematics and not yet usable as measurement: it needs
a renormalized, $0^{++}$-projected source in a fixed $t$-window, not more statistics.

\subsection{The exact missing step, and the ledger of this round}

Everything above converges on one step, and the wave's discipline requires that it be written as
sentences that can be decided rather than as an aspiration. There are two, and they are independent.

\begin{claimbox}{cM}
\claimhead{Missing step}{(two decidable sentences)} \emph{Along the physical line $\beta=2.5$,
uniformly in $(L,a)$ and inside the source's renormalization window $t\in[t_1,t_2]$ (with
$\sqrt{t_i}\sim c_i/m_G$):}
\begin{enumerate}
\item[(i)] \emph{``Quantitative isolation.''}
  \[
    s_2(t;a,L)\ \ge\ \frac{\bar E_t(a,L)-w_0\,s_1(t;a,L)}{1-w_0},
    \qquad
    \bar E_t=\tfrac12\bigl\langle\Omega|[o_t,[H,o_t]]|\Omega\bigr\rangle/\operatorname{Var}_\mu(o_t),
  \]
  \emph{with $\bar E_t$ computable and $s_1,s_2$ the two lowest support points of the
  $0^{++}$ measure. This is $I_2$ (isolation) in quantitative form; the threshold is fixed by a
  computable number, not by a free parameter.}
\item[(ii)] \emph{``Bottom-normalized single well.''} \emph{The effective potential in the binding
  coordinate must be a single well whose bottom is not pushed below the free value, with no Kramers
  doubling at the bottom. Equivalently}
  \[
    L_{0^{++}}\ \ge\ L_0=-\frac{d^2}{ds^2}
    \quad\text{(Dirichlet on }[0,L_s]\text{)}
    \qquad\Longrightarrow\qquad
    e_1\ \ge\ \frac{\pi^2}{L_s^2}.
  \]
\end{enumerate}
\emph{Given (i), the sharp inequality \eqref{eq:l2} yields $w_1\ge w_0$ immediately. Given (ii), the
rank-two trace \eqref{eq:ratio} yields $\rho^2\le C^2$ with $C=2$ \textbf{only} together with the
trace bound $T_2(V)\le5e_1$ of \code{research/v3/species/DT-DIRICHLET.md}: a one-sided floor supplies
a lower bound on $e_1$ and \emph{cannot} by itself imply $\rho^2\le C^2$ --- the witness is the local
profile with $e_1=18.0963\ge\pi^2$ and $\rho^2=27.1846$ --- so \textbf{the missing step is the ratio
bound itself}. Before both, none of the candidate routes of the bridge subsection moves.}
\end{claimbox}

\textbf{Note on the status of the two sentences.} Both are the same wall in another language. By the
campaign's own dictionaries the isolation-to-gap identification is fixed in the $\beta$/measure
coordinates, and the ratio coordinate's window is equivalent to a floor with a computable threshold,
\eqref{eq:eta}; so (i) and (ii) are $H1$ rewritten, not a species --- (i) is $H1$ seen through the
support point $s_2$, (ii) is $H1$ seen through a comparison operator. \textbf{This subsection is
therefore not a new species but a juxtaposition of three languages for one wall}, and the
identification is catalogued in \code{SIXTH-SPECIES-SYNTHESIS.md} \LC. The value of writing it this
way is diagnostic: whichever language the next round chooses, the object it must estimate is the
same, and only one of the three (the support point) has a sharp inequality already attached to it.

The second sentence is the one that carries the structural content, and it is worth spelling out why
the condition is two-sided. \emph{Exact} doubling at the bottom --- a symmetric doublet, a Kramers
pair, a broken discrete symmetry --- gives $\mu_1=\mu_2$ and hence $\rho^2=1$, which kills the lower
bound $\rho^2\ge1+c$ outright. \emph{Near}-degeneracy gives $\mu_1\to0$ with $\mu_2=O(1)$, hence
$\rho^2\to\infty$, which kills the upper bound $C^2$ just as outright. Both faces are the same
disease, and the discriminant is single: \textbf{is the bottom of the confining vacuum single?} That
is why the condition is stated as ``no Kramers doubling at the bottom'' rather than as a spectral
bound, and why it is a confinement statement (a single-valued vacuum, $Z_N$ center unbroken) rather
than a measure-theoretic one. Both sentences are decidable: (i) compares two numbers that are
either computable or defined, at a fixed inequality; (ii) is a multiplicity pattern of the lowest
few levels, which a lattice measurement can report without any new theory.

\textbf{The ledger.} The round's account, in the form the campaign keeps it:

\begin{center}
\footnotesize
\begin{tabular}{@{}>{\raggedright\arraybackslash}p{4.1cm}
  >{\raggedright\arraybackslash}p{7.6cm}c@{}}
\toprule
\textbf{Item} & \textbf{Statement} & \textbf{Label}\\
\midrule
Five species & HC, PS, TR, SW, LX closed, each with the number that closes it & \LN\\
Data-mode theorem & $d_N(\varepsilon)>0$ for every $N\le14$ \LE; minimum $0$ always \LT; the floor is the
unique discontinuity on the ray $\nu_h$, all continuous data moving as $0.183h$ \LT{} (values \LE); the
fit $d_N\sim e^{-0.4595N}$ is a fit on $N\le14$, not a theorem \LE & \LT/\LE\\
$(NA)$ & vacuous at $a_0=0$ (oddness), false uniformly in $a_0$ (sign flip
$-0.02307/+0.10750$), battlefield relocated to $\mathrm{SU}(N\ge3)$ & \LN\\
Counterexample & $G''=0.20511530-0.01105929\beta^2$ reproduced to five digits, $\beta_*=4.30661$ & \LE\\
Constant $C$ & $C=1.5$ refuted: five exact anchors violate $2.25$, and the physical band
$1.54$--$1.56$ gives $\rho^2\in[2.37,2.43]$, every point $>2.25$; $C=2$ is class-sharp only
(Dirichlet-type tower class), not structural --- ND box $\rho^2=9$, star $\rho^2=m+1$; floor
$\Delta/1.25\to\Delta/3$ & \LC\\
Radius & $1/48$ (SZZ, $N$-independent) \LT; $\lambda_{\mathrm{tH}}<1/24$ absent from the source
(bibliographic check, not a theorem) \LC & \LT/\LC\\
Bridge one & $\sum s_nw_n=\tfrac12 D_o/\operatorname{Var}(o)$, sharp $w_1\ge(\Lambda-\bar E)/(\Lambda-s_1)$,
threshold $s_2\ge(\bar E-w_0s_1)/(1-w_0)$ & \LT\\
Bridge two & Ky Fan $\mu_1+\mu_2=\min\operatorname{tr}(P_VLP_V)$, $\rho^2\le T_2(V)/\mu_1-1$,
$\rho^2\le C^2\iff e_1\ge t/(1+C^2)$ & \LT/\LN\\
$(TCI)$ & under-specified; the unique completion reduces to
$\lambda_{\min}(A_{\mathbb K}\mid\mathcal H_-)\ge cF+\mu$, and
$(\mathrm{TCI})\not\Rightarrow\Sigma\succeq\mu$; its only non-vacuous form is the condensate & \LN\\
Missing step & (i) quantitative isolation; (ii) bottom-normalized single well &
\LM\\
Clay & unbuilt & $[0\%]$\\
\bottomrule
\end{tabular}
\end{center}

\textbf{What this round generalized, and what it did not open.} The meta-theorem's emptiness
statement, ``computable $\wedge$ floor-sensitive $\Rightarrow\varnothing$'', was a qualitative claim
about four registered data modes. It is now a statement about the whole \emph{spectral class} ---
every finite list of weak-$*$ continuous functionals of the K\"all\'en--Lehmann measure --- and it
carries a number per mode, with an explicit defect profile for two of them. That is a genuine
widening: the four modes were the visible part of the class, and the class itself is now closed. The
immune map grows correspondingly, from thirty-three root families (Chapter~\ref{ch:immune}) to
\textbf{forty-four}: the five species, the ratio-coordinate family, and the alignment family are added,
each entry carries a label and a number, and the block-algebra entry is amended rather than added to
by the compensation finding of the $(TCI)$ subsection. The count is a bookkeeping convention, not a
theorem --- an independent re-attribution spans $45$--$47$ families --- and two load-bearing constants
are corrected.

\textbf{What it did not open: a route.} No new species was found. In the charter's taxonomy the round
produced $(\gamma)$ results only: five class-level no-gos, one quantified meta-theorem, one refuted
named inequality, one under-specified named inequality, two corrected constants, two bridges.
Options $(\alpha)$ (a fifth data mode) and $(\beta)$ (a new equivalent coordinate with a machine
strictly stronger than $H1$ somewhere) are untouched, and the chapter does not pretend otherwise.
Two further pieces of work are explicitly \emph{not} done and should not be reported as done. The
$\mathrm{SU}(N\ge3)$ version of $(NA)$ \emph{has} now been tested, and the earlier statement that it
was not is withdrawn: \code{research/v3/species/NA2-NONDEGENERATE.md} \S0 reports
$X(a_0=0)\neq0$ on a one-plaquette $\mathrm{SU}(3)$ model at magnitude $10^{-4}$ to
$1.7\times10^{-2}$, with the sign changing with the boundary holonomy, so $(NA)$ is refuted more
strongly than on $\mathrm{SU}(2)$; what stays open is the $4$D source $\partial_vI_3$, not the model.
Also open: the general statement
``$q\ge0\Rightarrow\rho^2\le4$'' has a numerical battery and a proof route but no proof (\LC). The
$\beta$-window of the provable $(TCI)$ floor below $\beta_c\approx0.76$ is likewise unmapped in $L$,
and is the cheapest open measurement the round leaves behind.

\textbf{What not to do next.} The negative results have an operational content, and the wave states
it so that the next round does not pay for it twice. Do not look for a ratio-specific mechanism:
\eqref{eq:ratio} and \eqref{eq:eta} show the ratio is a floor in disguise, the star graph
($\rho^2=m+1$ with diameter $2$ and $\mu_1=1$) shows no bound can depend only on local data, and the
double well shows no bound can depend only on upper-bound data. Do not retune $C$ towards $1.5$: the
anchors and the physical spectrum both refuse it. Do not spend compute on moment, angle, or
interlacing routes: their direction is fixed against the target, and the correction of the
interlacing candidate is quantitative ($\rho^2$ of a compressed space can land on either side:
$3.000<3.879<4.605$). What is worth doing is cheap and decisive: measure $\rho^2$ on a $(L,a)$ grid
(the third GEVP eigenvalue, the $0^{++}$ channel); measure $t=T_2(V)$ from two coarse operators and
compare it with $\mu_1$ through \eqref{eq:eta}; and report the multiplicity pattern of the lowest
$0^{++}$ levels, which is the direct test of sentence (ii).

\textbf{Closing.} The wall did not move. What moved is the accounting: five mechanism classes now
have receipts, the data-mode obstruction has a number attached to it, the last variance-free face
turned out to be $H1$ in disguise, the third named inequality turned out never to have been a
formula, and two constants that were silently wrong are now right. The gap remains exactly where it
was --- behind a conditioned variance in the infrared, a comparison operator in the ratio
coordinate, and a second support point in the alignment coordinate, which this round showed are the
same wall seen three times. \textbf{Clay: $0\%$.} No QED.
